Отсюда имеем:
\begin{equation}
	\tilde a \frac{\partial^2 \tilde u}{\partial \xi^2}+2\tilde b \frac{\partial^2 \tilde u}{\partial \xi\partial \eta}+
	\tilde c \frac{\partial^2 \tilde u}{\partial \eta^2} +
	\tilde{\Phi}(\xi,\eta,\tilde u, \frac{\partial\tilde u}{\partial \xi},\frac{\partial\tilde u}{\partial \eta}) = 0
	\tag{$\widetilde 1$}
\end{equation}
где:

\hspace*{2em} $\tilde a=a\xi^2_x+2b\xi_x\xi_y+c\xi^2_y$

\hspace*{2em} $\tilde b=a\xi_x\eta_x+b(\xi_x\eta_y+\xi_y\eta_x)+c\xi_y\eta_y$

\hspace*{2em} $\tilde c=a\eta^2_x+2b\eta_x\eta_y+c\eta^2_y$\\
Задача: выбрать такие $\xi$ и $\eta$, чтобы $\tilde a=\tilde c=0$, то есть:
$$ a\xi^2_x+2b\xi_x\xi_y+c\xi^2_y=0, \qquad (a\eta^2_x+2b\eta_x\eta_y+c\eta^2_y=0) $$
%ДИПСИЧНОЕ begin
\begin{note}
	Замена $(\xi,\eta)\leftrightarrow(x,y)$ предполагается невырожденной, т.е.
	\[
		J=\frac{\partial(\xi,\eta)}{\partial(x,y)} = \xi_x\eta_y - \xi_y\eta_x \neq 0.
	\]
	Это гарантирует существование обратного преобразования и эквивалентность уравнений.
\end{note}
%ДИПСИЧНОЕ begin
\begin{definition}[Характеристика и характерестическая кривая(случай $\mathbb{R}^2$, т.е. n=2)]
	Пусть функция $\omega(x,y) \in C^1(G):$ что на кривой $\omega(x,y)=0$ выполняется $\operatorname{grad}(\omega(x,y)) \neq \vec{0}$ и\\
	\begin{equation}
		a\omega^2_x+2b\omega_x\omega_y+c\omega^2_y=0
		\tag{2}
	\end{equation}
	Тогда кривая $\omega(x,y)=0$ называется характерестической кривой(характеристикой уравнения $(1)$), а само уравнение $(2)$ -- характерестическим уравнением уравнения $(1)$.
\end{definition}
\begin{lemma}[Связь первого интегралла и $\boldsymbol{\omega}$]%FIXME НАЗВАНИЕ КАКА
	Пусть $\omega(x,y) \in C^1(G)$ такова, что $\frac{\partial\omega}{\partial y}\ne 0$.\\
	Если $\omega(x,y)$ удовлетворяет уравнению $(2)$, то соотношение $\omega(x,y)=C$ является общим(первым) интеграллом ОДУ вида:
	\begin{equation}
		a\cdot\mathrm{d}y^2-2b\cdot\mathrm{d}x\mathrm{d}y+c\cdot\mathrm{d}x^2=0
		\tag{3}
	\end{equation}
	или $ a\cdot\mathrm{d}y^2-2b\cdot\mathrm{d}x\mathrm{d}y+c\cdot\mathrm{d}x^2=0 $, обозначим $(3')$\\
	Уравнение $(3)$ также называется характерестическим уравнением для $(1)$.
\end{lemma}
\begin{proof}
	Т.к. $\omega(x,y)\in C^1(G),\,\omega_y\ne 0$ и $\omega(x,y)$ удовлетворяет уравнению $(2)$, то из него получаем:\\
	$\displaystyle(2'):\; a\left(\frac{\omega_x}{\omega_y}\right)^2+2b\left(\frac{\omega_x}{\omega_y}\right)+c=0$ или
	$\displaystyle a\left(-\frac{\omega_x}{\omega_y}\right)^2-2b\left(-\frac{\omega_x}{\omega_y}\right)+c=0$
\end{proof}
Вспомним, что соотношение $\omega(x,y)=C$ является общим интегралом ОДУ, если это соотношение ($\omega(x,y)=C$) неявно задает функцию $y=y(x)$, которая
является решением указанного ОДУ.\\
Тогда по теореме ``О неявных функциях'', из $\omega_y\ne0\rightarrow$

\hspace*{1em}  1. Гиперболический тип: пусть в некоторой области $G$, $a\ge0$, $D>0$, тогда уравнения $(3)$ и $(3')$ -- различны и у них разлиные незвисимые 1\textsuperscript{ые} интеграллы: $\phi(x,y)=C_1$, $\psi(X,y)=C_2$ и $\displaystyle \frac{\partial(\phi,\psi)}{\partial(x,y)}\ne 0$\\%FIXME ЧТО ТУТ ПРОИСХОДИТ? ПЕРЕМЕШААААААЛИ??? короче структуру надо поправить на основе других лекций
Перейдем в уравнении $(1)$ к независимым переменным:
$\xi=\phi(x,y),\,\eta=\psi(x,y)$.\\
Тогда уравнение $(1)$ в некоторой окрестности из $G$ примет вид:

\hspace*{2em} $2\widetilde b \widetilde u_{\xi\eta}+\widetilde \Phi=0$, откуда $\widetilde u_{\xi\eta}=\widetilde{\widetilde \Phi}$ --- 1\textsuperscript{ая} каноническая форма.

\hspace*{1em}  2. Эллиптический тип: $\displaystyle \hat{u}_{\alpha\alpha}+\hat{u}_{\beta\beta}=\hat\Phi $;

\hspace*{1em}  3. Параболический тип: $\displaystyle \widetilde{u}_{\eta\eta}=\widetilde{\widetilde{\Phi}} $;
\begin{note}
	Если в некоторой области $G$ уравнение имеет гиперболический тип(ГТ), то в этой области $G$ существует 2 семейства $\mathbb R$(действительных) - характеристик, а если эллиптический тип, то в этой области $G$ существует 2 семейства $\mathbb C$(комплексно) - сопряженных характеристик.
\end{note}

